
\section{研究背景与动机}
物理模拟的探索源于20世纪中叶有限元方法在工程领域的应用 [1]。自20世纪90年代以来，计算机辅助工程（CAE）与计算机图形学（CG）从不同维度推动了物理模拟的发展。CAE 侧重于通过高精度数值方法实现工业设计的可靠预测 [2]；而 CG 则强调在物理机制的基础上提升视觉表现的真实感与交互效率 [3]。尽管侧重点各异，二者在算法核心上均追求数值模拟的精确性与稳定性。

从计算机科学视角看，物理模拟是计算几何、数值分析与高性能计算等学科的交叉产物。研究者通过算法创新，在有限算力下实现了跨尺度的仿真任务。随着人工智能技术的发展，高保真仿真环境在具身智能等前沿应用中日益重要 [4]。由于智能体的训练涉及大量闭环迭代，任何数值不稳定性都会严重影响学习效率。因此，攻克仿真中的数值瓶颈已成为提升算法鲁棒性与计算效率的核心要求。

在深入探讨具体算法之前，必须审视物理模拟在数值计算层面的核心逻辑。以典型的动力学问题为例，其核心在于通过空间离散化与代数表达，将连续的偏微分方程转化为代数系统：
\begin{equation}
    \mathbf{M}\mathbf{\ddot{x}} + \mathbf{f}_{int}(\mathbf{x}, \mathbf{\dot{x}}) = \mathbf{f}_{ext}
\end{equation}
其中全局系统的组装过程不仅涉及复杂的几何拓扑处理，更直接决定了后续求解器的计算效率与鲁棒性。

当前物理仿真的核心瓶颈在于系统矩阵的高条件数，即系统的病态性 [5]。当系统中存在物理属性或数值刚度极度失配的“软硬耦合”成分时，系统特征值的频谱分离会导致条件数激增。这种现象在不可拉伸细杆模拟（极高拉伸刚度与微弱弯曲刚度）、有限元方法（局部畸变产生的人造刚度）以及模态综合法（高频强刚性与全局柔顺振动）中均有体现。

系统病态性会显著影响仿真效率与稳定性。在底层求解阶段，高条件数会降低克雷洛夫子空间方法的收敛速度并累积舍入误差 [6]。在时间积分方面，病态性迫使半隐式方法采用极小步长，或导致隐式方法中的牛顿迭代难以收敛。在图形学应用中，这些数值问题常表现为物体的运动迟滞、能量震荡或拓扑交错等视觉瑕疵 [7]。

为了攻克系统病态性带来的挑战，传统思路往往侧重于在已有离散表示上施加代数层面的预处理技术。这类方法作为通用求解工具，通过在代数层面构造近似逆算子，有效提升了大规模线性系统的收敛速度。然而，在本文关注的若干极端病态场景中（如极高的弹性模量比或恶劣的几何退化），困难并不主要来自系统规模本身，而是来自局部高频硬成分与全局低频软响应并存所导致的条件数恶化。由于纯代数方法较难从物理源头感知并切断这些高频刚度，其预处理效能往往会随着硬约束的增强而显著衰减。因此，本文主张回归计算机科学中“表达决定性能”的本质，从离散自由度与基函数的设计出发，研究如何在表示空间层面隔离局部高频刚度，并自然保留全局低频响应。

相比于代数预处理，这种“物理感知（Physics-aware）”的重构方法具有更高的计算效率与更强的算法防御性。本研究将这一核心思想应用于解决三个典型的软硬耦合难题：
\begin{enumerate}
    \item 不可拉伸弹性细杆仿真：针对毛发与纱线级布料动画，通过基函数重构消除高纵向刚度带来的数值限制。
    \item 畸变单元稳定性处理：解决非结构化网格中局部几何畸变导致的系统矩阵崩溃现象 [8]。
    \item 广义特征值问题优化：针对图形学模态分析，通过更稳定的基空间表征提升特征谱提取的精度。
通过这三个问题，本文展示了硬约束的抽象程度逐渐增加的递进关系——从显式的空间距离，到隐蔽的局部拓扑约束，最终抽象为频域中的相位缺失。
\end{enumerate}

通过回顾并整合这一系列工作，本论文旨在引出一种全新的研究范式——即通过面向软硬耦合问题的基函数设计，构建更健壮、更高效的物理仿真系统。为了清晰地阐述这一范式，本章接下来的内容将首先建立物理模拟的基础离散管线视角，随后对比纯代数预处理与表示空间重构两条技术路线的本质差异，并最终引出本文具体关注的三类典型病态系统及其面临的数值困境。